\documentclass[conference]{IEEEtran}
\usepackage{times}
\usepackage[numbers]{natbib}
\usepackage{multicol}
\usepackage{graphicx}
\usepackage{amsmath,amssymb}
\usepackage[bookmarks=true]{hyperref}

\newcommand{\ourhand}{CoDe Hand}
\title{Hand Design}
\author{Author Names Omitted for Anonymous Review}
\hypersetup{pdfauthor={},pdftitle={CoDe Hand: Same Hand, Different Contact}}

% RSS 2027 Stage 1: at most 5 pages of content + 1 page of references.
% Section files follow the PLATO-RAL project structure.
\begin{document}
\maketitle
% ---- sections/00_abstract.tex ----
\begin{abstract}
% TODO: Motivation, integration problem, design, and principal evidence.
% Write this last; do not state planned experiments as completed results.
The abstract goes here.
\end{abstract}

% ---- sections/01_introduction.tex ----
\section{Introduction}
\label{sec:introduction}
% Motivation: an accessible robot-learning platform for different contact interfaces.
% Problem: module integration can require hand redesign.
% Research question, scope, and evidence-supported contributions.

% ---- sections/02_related_works.tex ----
\section{Related Work}
\label{sec:related_work}
% Accessible dexterous hands for robot learning.
% Contact interfaces, tactile sensors, and their integration requirements.
% Cite only literature directly tied to the central problem and novelty.

% ---- sections/03_design.tex ----
\section{Design}
\label{sec:design}
% Design requirements and supported module envelope.
% Servo-based finger architecture and crossed four-bar mechanism.
% Proximal/distal modules, mounting, replacement, and model updates.
% Add dimensions and hardware specifications only when verified.

% ---- sections/04_experiments.tex ----
\section{Experiment}
\label{sec:experiment}
% Hypotheses, baselines, methods, and metrics.
% Module accommodation and usable motion versus a fair serial-2R baseline.
% Supported-module exchange versus first-time module integration cost.
% Learning-policy execution with functioning contact/sensing modules.
% Distinguish CAD, simulation, and physical evidence; report negative results.

% ---- sections/05_conclusion.tex ----
\section{Conclusion}
\label{sec:conclusion}
% Supported findings, limitations, and remaining validation.


\maketitle


\bibliographystyle{plainnat}
% Enable after adding the first verified citation to the manuscript.
% \bibliography{references}
\end{document}
